\subsection{Searching for \VCs}
\label{sec:guided-search}

\begin{algorithm}[tb]
\footnotesize
\caption{Parallel Beam Search}\label{alg:beam-algo}
\SetKwProg{Fn}{Function}{}{end}
\KwData{All causal relationships between faults $\mathbb E$}
\KwData{Beam size B}
\KwResult{All \vcs $\mathbb C$}
\BlankLine
\Fn{beamSearch($\mathbb E$, B)}{
    queue $\gets e \in \mathbb E$\;
    $\mathbb C \gets \emptyset$\;
    \While{queue $\neq \emptyset $}{
        next $\gets \emptyset$\;
        \ForPar{$c \in \mathrm{queue}$}
        {
            \ForPar{$e \in \mathbb E$}
            {
                \If{match(c.lastEdge, e)}
                {
                    new $\gets$ c.append(e)\;
                    \eIf{isCycle(new)}
                    {
                        $\mathbb C$.add(new)   
                    }
                    {
                        next.add(new)
                    }
                }
            }
        }
        sort(next)\;
        queue $\gets$ next.subList(0, B)\;
    }
    \Return{$\mathbb C$}\;
}
\Fn{match(edge1, edge2)}{
\Return{edge1.interference == edge2.injectedFault \& isCompatible(edge1.state, edge2.state)}\;
}
\Fn{isCycle(c)}{
\Return{match(c.lastEdge, c.firstEdge)}\;
}
\end{algorithm}


Due to the large search space 
and the need to 
incorporate the local compatibility check, 
\tool uses a customized \textbf{parallel beam search} 
to find cycles, 
as shown in \Cref{alg:beam-algo}. 
The algorithm starts from all causal relationships obtained from the three-phased fault injection, forming propagation chains of length 1 (line 2).
Each while-loop iteration (line 4) represents one-level on the search tree, where one edge \texttt{e} is connected to each propagation chain \texttt{c} in the \texttt{queue} (lines 6-13). Before connection, \tool performs the compatibility check (line 17) between the last edge in the chain and the new edge (line 8). If each chain cycles back to its beginning (line 10), \tool reports a potential \vc.




At each level of the beam search, only $B$ active chains are kept, sorted by the average intra-cluster interference similarity score (\Cref{sec:alloc:3pa}) of the injected faults in the chain (lines 14-15). The intuition is that \tool favors \vcs involving complex error handling logic, which developers might overlook during testing. Suppose each chain $C_i$ has $r$ fault injections $f_{k_1}, f_{k_2}, ..., f_{k_r}$, and each fault $f_{k_j}$ belongs to fault cluster $G_j$. The score used in ranking is defined as: 
$ \mathrm{Score}(C_i) = \Sigma_{j=0}^r \mathrm{SimScore}(G_j) / \lVert C_i \rVert$. 
Chains with lower intra-cluster similarity are kept, as they likely involve conditional causal relationships.

\tool perform the beam search on the entire fault causal space explored by the fault injection. 
Although there is no limit on the number of faults in each chain, chains cannot grow indefinitely as the fault injection runs in the chain must remain compatible.


\textbf{Clustering Reported Cycles.} 
Because in 3PA's phase two and three, 
\tool randomly picks a fault in a cluster of \textit{causally equivalent} faults 
to perform injection, 
our detection algorithm can identify 
equivalent \vcs 
which contains \textit{causally equivalent} faults. 
For example, suppose \tool reports two cycles ``$f_1 \rightarrow f_2 \rightarrow f_1$'' and ``$f_3 \rightarrow f_2 \rightarrow f_3$'', while $f_1$ and $f_3$ are from the same fault cluster $G$, the two cycles reported are likely the same bug due to $f_1$ and $f_3$'s similar interference on the system.
\tool automatically clusters cycles found in the beam search based on the fault clusters (\Cref{sec:alloc:3pa}) involved in the cycle. 
